\documentclass[11pt]{article}
\usepackage[T1]{fontenc}
\usepackage{lmodern,microtype}
\usepackage{amsmath,amssymb,booktabs,enumitem}
\usepackage[letterpaper,margin=1in]{geometry}
\usepackage[colorlinks=true,urlcolor=blue,citecolor=blue]{hyperref}
\newcommand{\Q}{\mathbb Q}
\newcommand{\F}{\mathbb F}
\newcommand{\Z}{\mathbb Z}
\title{Metaplectic Jacquet--Langlands Transfer, Companion Paper (Rank Two, Wild Place)\\Revision Log: Draft 4 to Draft 5}
\author{Prepared by Claude (Anthropic) for David Joyner}
\date{November 8, 2026}
\begin{document}
\maketitle
\section*{Scope}
Draft 5 (i) incorporates the two clarifications requested in ChatGPT's
referee report on Draft 4 (recommendation: minor revision), and (ii) adds
one new theorem: Target B holds for every irreducible representation of
$D^\times$ trivial on $K^3$.  Everything else is unchanged.

\section*{The two clarifications}
\begin{enumerate}[leftmargin=*]
\item Lemma \texttt{lem:centralizers}(a) is now stated for $2$-elements
$y$ whose image in $D_8$ is nontrivial, with a parenthetical explaining
that the central $\mu_2\subset T$ of the odd case (whose centralizer is
all of $\Gamma_0$) is excluded; the Brauer proof of Theorem
\texttt{thm:reduction-all} only invokes the lemma for such $y$.
\item Lemma \texttt{lem:H1consequences}(c): the splitting
$\tilde V\simeq\mu_3\times V$ is now derived from $\tilde V$ being
abelian \emph{of exponent $3$} (cubes of lifts of $K^3$ are star lifts
of elements of $K^9$, on which $\psi_6=1$; cubes of central lifts are
star lifts of cubes, on which $\psi_6=\theta_0^3=1$), with the remark
that abelianness alone would not suffice ($C_9\to C_3$).
\end{enumerate}

\section*{Target B at level two (new Section 12.4)}
For $\rho$ irreducible, trivial on $K^3$, nonscalar on $K^2$:
$\rho|_{K^1}=\chi_3+\chi_3'$ (the two $3$-dimensional characters of the
Heisenberg group $Q_1/Q_3$), $\dim\rho=6$, four such $\rho$ per
normalized central character.  Findings, all by exact computation
(scripts 13--15, about ten minutes):
\begin{enumerate}[leftmargin=*]
\item The sufficient criterion of Proposition \texttt{prop:criterion}
\emph{fails} for these $\rho$: on the $219$ good classes of $H=Q_3/Q_7$,
$F_\rho$ is a rational but not integral combination of the
$\mathrm{Ad}(Q_1)$-invariant virtual characters of $H$ (Smith normal
form over $\Z$).  For the three orbit sums of linear characters of the
Heisenberg group it succeeds, consistent with Theorem \texttt{thm:levelone}.
\item Nevertheless $\Phi_\rho|_{\tilde U_3}$ is a genuine character.
Since $\Phi_1|_{\tilde U_3}=\lambda_0\cdot D$ with $D=\sum_{h\in\Lambda}\zeta^h$,
$\Phi_\rho|_{\tilde U_3}=\lambda_0\cdot(\Theta_\rho\circ r)D$, and the
exact character table of $H$ ($297$ classes, $243$ linear and $54$
three-dimensional characters) gives
$((\chi_3+\chi_3')\circ r)\cdot D=\sum\{\chi'\in\mathrm{Irr}(H):\chi'(1)=3\}$,
each of the $54$ exactly once.
\item Clifford theory then reduces Target B for $\rho$ to the freeness of
the $\mathrm{Ad}(Q_1)$-action on the $54$ genuine three-dimensional
irreducibles $\lambda_0\chi'$ of $\hat H=\tilde U_3/\ell(K^7)\ker\tilde\theta$
(a size-$9$ orbit would give a non-integral multiplicity $e|O|/27$).
Computing $\lambda_0$ on $\hat H$ in the cocycle model and the twisted
action $\mu^t=\lambda_0 c_t(\chi'\circ\mathrm{Ad}\,t)$: exactly two
orbits, both of size $27$.  Hence
\[
 \tau_\rho|_{\tilde U_1}=(\varepsilon_\rho\circ p)\otimes\bigl(\mathrm{Ind}(\lambda_0\otimes\chi_1')\oplus\mathrm{Ind}(\lambda_0\otimes\chi_2')\bigr),
\]
two irreducibles of dimension $81$, $\dim\tau_\rho=162$ (Theorem
\texttt{thm:leveltwo}, [P$^c$]).  Cross-checks: the untwisted action on
the same $54$ characters has six orbits of size $9$ (so the
$\lambda_0$-twist is essential), and the $243$ genuine linear characters
form nine free orbits, recovering Theorem \texttt{thm:levelone}.
\end{enumerate}

\section*{The cube-root correspondence (new Section 12.5, [O])}
Both level theorems fit one pattern: to each irreducible $\chi$ of
$Q_1/Q_{m+1}$ corresponds an irreducible $\chi'$ of $Q_3/Q_{3m+1}$ of the
same degree with $\Phi_1\cdot(\chi\circ r)=\Theta_{\mathrm{Ind}(\lambda_0\otimes\chi')}$;
this is formulated as a conjecture implying Target B for all $\rho$,
and reduced to two statements: (i) $\mathrm{Ad}(Q_1)$ acts freely on the
genuine irreducibles of $\hat H_m$ (equivalently every
$\mathrm{Ind}_{\tilde U_3}^{\tilde U_1}\mu$ is irreducible), a statement
about the cover alone; (ii) $(\chi\circ r)D$ is a virtual character of
$Q_3/Q_{3m+1}$.  Both hold for $m\le2$.  For $m=3$ the group has order
$3^{11}$; Burnside's algorithm as implemented will not reach it, but
Clifford theory relative to $Q_4/Q_{10}$ would.

\section*{Computational notes}
\texttt{step13\_level2.py} builds $H$ from the Draft 1 model of $Q/Q_7$
(\texttt{finite.pkl}), computes its character table numerically
(Burnside) and re-verifies it exactly in $\Z[\zeta_3]$ (all
orthogonality relations), and solves the criterion over $\Z$ with a
hand-written Smith normal form (\texttt{zsolve.py}).
\texttt{step14\_level2\_restriction.py} decomposes $(\chi\circ r)D$.
\texttt{step15\_level2\_clifford.py} computes $\lambda_0$ on $\hat H$
(formula \eqref{eq:lambda0} of the paper, as in script 11), the
conjugation action with the exact cocycle, and the orbits; it also
checks on $300$ random pairs that $\lambda_0$ descends to a character of
$\hat H$ (a consistency check, not load bearing).  A scalar-extraction
bug in a first version of script 15 (wrong coordinate for the
$F$-component) was caught by that check.

\section*{Status}
Theorem \texttt{thm:leveltwo} [P$^c$] (new); Theorems
\texttt{thm:reduction-all}, \texttt{thm:H1-gen} [P] (clarifications
applied); Section 12.5 [O]; nothing else changed.
\end{document}
